\section{From ideal polyhedra to the Euclidean limit of WP geometry}
\label{sec:flatlimit}

\subsection{A prescribed intrinsic interpolation}

For the weights in~\eqref{eq:flatweights}, define
\begin{equation}
 \theta_i(t)=2\pi t(1-\mu_i),\qquad
 a_i(t)=1-t(1-\mu_i),\qquad 0\le t<1.
 \label{eq:anglepath}
\end{equation}
Hold the marked conformal sphere $X\in\M_{0,n}$ fixed, and let $g_t(X)$
be its curvature-$-1$ cone metric with these openings. Its area is
\begin{equation}
 A_t=2\pi\left(\sum_i a_i(t)-2\right)
    =2\pi(n-2)(1-t).
 \label{eq:areapath}
\end{equation}
At $t=0$ the metric is the complete cusped metric, realized by an ideal
polyhedron, including doubled ideal polygons. For $0<t<1$ it is a positive-deficit hyperbolic cone metric,
realized by a compact convex hyperbolic polyhedron, with flattened
realizations included when necessary. The realization results are
reviewed in~\cite{Fillastre}; the ideal and generalized settings are
treated in~\cite{RivinIntrinsic,Schlenker}.

Normalize to unit area:
\begin{equation}
 \widehat g_t=A_t^{-1}g_t,
 \qquad \operatorname{Area}(\widehat g_t)=1,
 \qquad K(\widehat g_t)=-A_t.
 \label{eq:unitareanormalization}
\end{equation}
The unit-area uniformization and local deformation results of
Mazzeo--Weiss give the flat limit
\begin{equation}
 \widehat g_t(X)\longrightarrow g_{\mathrm E,\mu}(X)
 \quad (t\to1),
 \label{eq:flatmetriclimit}
\end{equation}
with the Euclidean metric normalized to area one. Here their angle
parameters are $\beta_i(t)=-a_i(t)\in(-1,0)$ near $t=1$, and
$\chi(S^2)+\sum_i\beta_i(t)=-D(1-t)\le0$, where $D=n-2$.
Thus their nonpositive-curvature existence and uniqueness theorem
applies, including the flat endpoint. Their smooth local slices give
convergence away from the cone points and uniform conic estimates
after a local trivialization fixing the marks
\cite[Theorem 1, Section 6.1, and Proposition 16]{MazzeoWeiss}.
This is a deduction along the specified path; it is local over
moduli where the marks remain separated.

This constructs a family of complex fixed-angle leaves with two
geometric endpoints: cusped hyperbolic surfaces/ideal polyhedra, and
flat cone spheres/Euclidean shapes. It does not assign a single
complex structure to all angle variables in the full real space of
hyperbolic polyhedra. Nor is it generally radial scaling of a fixed
Euclidean vertex configuration in the Klein ball. Such radial scaling
makes all vertices ideal simultaneously only for a cospherical
configuration. Every nondegenerate Euclidean tetrahedron is
cospherical, but its radial path still need not preserve $X$ or
the angle path~\eqref{eq:anglepath}.

\subsection{The normalized local WP form}

The scale of the WP form can be seen directly from the cotangent
pairing. Write $w=x+iy$ and $g_t=\rho_t|\dd w|^2$. For
$q=\varphi\,\dd w^2$ and $r=\psi\,\dd w^2$, the raw coefficient
pairing is
\[
 Q_t(q,r)=\int_X\frac{\varphi\overline\psi}{\rho_t}\,\dd x\,\dd y.
\]
Since $\rho_t=A_t\widehat\rho_t$, the rescaled cometric $A_tQ_t$
has a finite flat limit. A quadratic differential has at most a
simple pole at a marked point, while
$\widehat\rho_t\to\rho_{\mathrm E}$ with
$\rho_{\mathrm E}=c_i(w)|w|^{-2\mu_i}$ for a smooth positive local
factor $c_i$. The limiting integrand is
$O(|w|^{2\mu_i-2})\,\dd x\,\dd y$, which is integrable because
$\mu_i>0$. On a compact subset of a coordinate chart in moduli, the
conic estimates bound the nearby integrands by
$C|w|^{2m-2}\,\dd x\,\dd y$ for some $m>0$. Dominated convergence
therefore gives locally uniform convergence of the pairing matrices.
Their flat limit is positive definite, so matrix inversion gives a
tangent metric and a K\"ahler form of order $A_t$. No assertion here
allows the marked points to collide or a weight to tend to zero.

To determine the constant, use the raw tangent coefficients
$G_{j\bar k}(g)=\int_X\nu_j\overline{\nu_k}\,\dd A_g$, where
$\nu_j$ is the harmonic Beltrami representative of a complex coordinate
vector $\partial/\partial z_j$. Our FN normalization is
$\omega_{\WP}=i\sum G_{j\bar k}\,\dd z_j\wedge\dd\bar z_k$,
as fixed in Section~\ref{sec:conventions}; equivalently, on real
tangent vectors it is $-2\operatorname{Im}G$
\cite[Section 6.1, equation (38)]{KawazumiWP}.
Wang's period-variation calculation, translated to the positive area
pairing~\eqref{eq:DMarea}, gives
\[
 -\frac{\partial^2\log\mathcal A}{\partial z_j\partial\bar z_k}
  =\frac{G_{j\bar k}(g_{\mathrm E})}
         {\operatorname{Area}(g_{\mathrm E})}.
\]
Indeed, the component of the variation of $[\eta]$ orthogonal to
$[\eta]$ is represented by $\nu_j\mathbin{\lrcorner}\eta$; its
negative Hermitian norm divided by $H(\eta,\eta)$ is the displayed
area average. This is the normalization-independent content of
\cite[Theorem 1.2.1 and Section 4.4]{WangFlatWP}. Inserting
\eqref{eq:CHpotential} now gives
\begin{equation}
 \boxed{\quad A_t^{-1}\omega_{\WP,t}
       \longrightarrow\frac12\omega_{\mathrm{CH},\mu}.\quad}
 \label{eq:localWPCHlimit}
\end{equation}
The convergence is of form coefficients locally uniformly on the
open moduli space. The factor is also checked independently by the
class calculation below. This conclusion combines the two cited
inputs with the specified FN convention; neither source states this
hyperbolic-to-flat limit verbatim. Local convergence on a noncompact
moduli space alone does not justify convergence of total volumes.

\subsection{The class calculation and an exact final-chamber polynomial}

Put $d=n-3$, $D=n-2$, and $\mu_I=\sum_{i\in I}\mu_i$. Let
$\delta_I=\delta_{I^c}$ denote the divisor of the stable bipartition
$I\mid I^c$. Define the real divisor class
\begin{equation}
 B_\mu=\frac12\left(
   -\sum_i\mu_i\psi_i+
   \sum_{\{I,I^c\}}\min(\mu_I,\mu_{I^c})\delta_I\right),
 \label{eq:CHclassB}
\end{equation}
where the second sum counts each unordered stable bipartition once.
Use the real compactification class on the ordinary
Deligne--Mumford compactification, defined at balanced weights by
continuity. Its normalization and volume formula are
\begin{equation}
 [\omega_{\mathrm{CH},\mu}]=4\pi B_\mu,\qquad
 V_{\mathrm{CH}}(\mu)=\frac{(4\pi)^d}{d!}
                     \int_{\MM_{0,n}}B_\mu^d.
 \label{eq:CHintersection}
\end{equation}
Away from balanced weights, Koziarz--Nguyen's normalized form
$\Omega_\mu$ represents $B_\mu$ as a closed current. Their
Section 7.3 gives the divisor class, and Section 7.4 converts it
to curvature $-1$ by
$\omega_{\mathrm{CH},\mu}=4\pi\Omega_\mu$
\cite[Theorem 1.1 and Sections 7.3--7.4]{KoziarzNguyen}.
Their direct metric/current analysis excludes subsets $\mu_I=1$;
at these balanced weights the displayed divisor class and volume
formula are understood by continuity, as explained in their
Remark 1.3. The volume on the left is the finite volume of the
open labeled moduli space, with no permutation quotient.

Let $\varepsilon=1-t$, $r_i=1-\mu_i$, and
$r_I=\sum_{i\in I}r_i=|I|-\mu_I$. Then
$a_i=\mu_i+\varepsilon r_i$ and $A=2\pi D\varepsilon$. Write
$\mathcal L_\mu$ for the stable subsets with $\mu_I<1$, each light
side counted once. For sufficiently small positive $\varepsilon$,
these are precisely the light subsets in the adjacent hyperbolic
chamber. Subsets with $\mu_I=1$ are \emph{not} light along this path.
Define
\begin{equation}
 \boxed{\quad
 \Xi_\mu=\frac{1}{2D^2}\left(
 -\sum_i(1-\mu_i)^2\psi_i+
 \sum_{I\in\mathcal L_\mu}(|I|-\mu_I)^2\delta_I\right).
 \quad}
 \label{eq:Xiclass}
\end{equation}

\begin{proposition}[A deduction from the two compactification formulas]
In the final chamber of the path~\eqref{eq:anglepath},
\begin{align}
 [\omega_{\WP}(A)]
   &=\frac A2[\omega_{\mathrm{CH},\mu}]+A^2\Xi_\mu,
       \label{eq:exactWPclassflat}\\
 V_{\WP}(A)
   &=\frac{A^d}{d!}\int_{\MM_{0,n}}
       \left(\frac12[\omega_{\mathrm{CH},\mu}]+A\Xi_\mu\right)^d.
       \label{eq:exactWPvolumeflat}
\end{align}
In particular,
\begin{equation}
 \boxed{\quad
 V_{\WP}(A)=2^{-d}A^d V_{\mathrm{CH}}(\mu)+O(A^{d+1}).
 \quad}
 \label{eq:WPCHvolumeasymptotic}
\end{equation}
The volume in~\eqref{eq:exactWPvolumeflat} is exactly a polynomial in
$A$ with possible degrees $d,d+1,\ldots,2d$ within that chamber.
\end{proposition}

\begin{proof}
In terms of deficit weights, the pulled-back WP class
\eqref{eq:wallclass}, divided by $2\pi^2$, is
\[
 \beta(a)=\kappa_1-\sum_i(1-a_i)^2\psi_i+
           \sum_{a_I<1}(1-a_I)^2\delta_I.
\]
Here all collision sums run over stable subsets, so
$2\le |I|\le n-2$; since $\sum a_i>2$, both sides of a
bipartition cannot be light. We verify the two needed endpoint
identities on the ordinary universal curve
$\pi:\mathcal C\to\MM_{0,n}$. Write $\Gamma_i$ for its sections,
$K=K_{\mathcal C/\MM_{0,n}}+\sum_i\Gamma_i$, and $F_I$ for
the vertical component over $\delta_I$ carrying the marks in $I$.
For the fixed final-chamber list $\mathcal L_\mu$, set
\[
 \Lambda(a)=K-\sum_i(1-a_i)\Gamma_i
                  -\sum_{I\in\mathcal L_\mu}(1-a_I)F_I.
\]
The universal-curve intersections used in
\cite[Section 8.4]{KoziarzNguyen} are
\[
 \begin{gathered}
 \pi_*K^2=\kappa_1,\qquad K\Gamma_i=0,\qquad
 \pi_*(\Gamma_i^2)=-\psi_i,\qquad \Gamma_i\Gamma_j=0\ (i\ne j),\\
 \pi_*(KF_I)=(|I|-1)\delta_I,\qquad
 \pi_*(F_IF_J)=-\mathbf1_{I=J}\delta_I,\qquad
 \pi_*(\Gamma_iF_I)=\mathbf1_{i\in I}\delta_I.
 \end{gathered}
\]
Consequently $\pi_*\Lambda(a)^2=\beta(a)$. For example, with
$c=1-a_I$, the coefficient of $\delta_I$ in this square is
$-2c(|I|-1)+2c\sum_{i\in I}(1-a_i)-c^2=c^2$; this checks
the sign of the vertical correction.

For rational flat weights away from balanced subsets, the corrected
divisor $\Lambda(\mu)$ is fiberwise trivial and is the pullback of
Koziarz--Nguyen's extended period-line class $B_\mu$
\cite[Sections 8.1 and 8.4]{KoziarzNguyen}. All coefficients extend
continuously to real weights and to balanced subsets, whose vertical
coefficient $1-\mu_I$ vanishes. Hence
$\Lambda(\mu)=\pi^*B_\mu$ as real divisor classes throughout
\eqref{eq:flatweights}. Along our fixed chamber,
\[
 \Lambda(\mu+\varepsilon r)=\pi^*B_\mu+\varepsilon R,
 \qquad R=\sum_i r_i\Gamma_i+\sum_{I\in\mathcal L_\mu}r_IF_I,
 \qquad \pi_*R=\sum_i r_i=D.
\]
The projection formula now gives
$\pi_*\Lambda(\mu)^2=0$ and first-order term
$2\varepsilon D B_\mu$. Equivalently,
\begin{align*}
 \beta(\mu)&=0,\\
 \sum_i r_i^2\psi_i-
   \sum_{I\in\mathcal L_\mu}(1-\mu_I)r_I\delta_I&=D B_\mu.
\end{align*}
These are general class identities obtained by pushforward; finite
intersection checks are supplementary.

Expanding the quadratic expression for $\beta$ now gives
\[
 \beta(\mu+\varepsilon r)=2\varepsilon D B_\mu+
 \varepsilon^2\left(-\sum_i r_i^2\psi_i+
                \sum_{I\in\mathcal L_\mu}r_I^2\delta_I\right).
\]
Multiplication by $2\pi^2$, followed by
$A=2\pi D\varepsilon$ and~\eqref{eq:CHintersection}, proves
\eqref{eq:exactWPclassflat}. The WP intersection-volume theorem
\cite{AMN} proves~\eqref{eq:exactWPvolumeflat}; its leading term is
\eqref{eq:WPCHvolumeasymptotic}.
\end{proof}

Within the original parameter interval $0<\varepsilon\le1$, the
open final chamber is
\begin{equation}
 0<\varepsilon<\varepsilon_*,\qquad
 \varepsilon_*=
 \begin{cases}
 \displaystyle\min_{I\in\mathcal L_\mu}
                \frac{1-\mu_I}{|I|-\mu_I},&\mathcal L_\mu\ne\varnothing,\\
 1,&\mathcal L_\mu=\varnothing.
 \end{cases}
 \label{eq:finalchamberwidth}
\end{equation}
When the light-subset set is empty, the same polynomial also holds
at the cusped endpoint $\varepsilon=1$. At a collision wall the
vanishing correction makes the adjacent volume polynomials agree,
but the strict inequality specifies which chamber supplies the
coefficients. Along the full
path, a stable subset crosses its wall at
$t_I=(|I|-1)/(|I|-\mu_I)$ when $\mu_I<1$. Thus the ideal-to-flat
volume is piecewise polynomial, even though the final expansion is
an exact polynomial.

\subsection{Examples and nonuniformity at balanced weights}

With four equal weights $\mu_i=1/2$, all openings are $\pi t$, and
there are no active collision corrections for $t<1$. Thus
\begin{equation}
 V_{\WP}(t)=2\pi^2(1-t^2),\qquad A=4\pi(1-t),\qquad
 V_{\WP}(A)=\pi A-\frac{A^2}{8}.
 \label{eq:equalflatvolume}
\end{equation}
Here $\int B_\mu=1/2$, $V_{\mathrm{CH}}=2\pi$, and
$\int\Xi_\mu=-1/8$, checking all constants directly.

For $\mu=(3/10,2/5,3/5,7/10)$, the light pairs are $\{1,2\}$
and $\{1,3\}$, while $\{1,4\}\mid\{2,3\}$ is balanced.
One obtains $\int B_\mu=3/10$ and $\int\Xi_\mu=9/40$.
The two positive wall parameters are $\varepsilon=3/13$ and $1/11$,
so the final-chamber volume is
\[
 V_{\WP}(A)=\frac35\pi A+\frac9{40}A^2,
 \quad 0<A<\frac{4\pi}{11},\qquad
 V_{\mathrm{CH}}=\frac65\pi.
\]
At the flat endpoint the uncorrected cone polynomial would equal
$-\pi^2/5$; the two light-pair corrections, $9\pi^2/50$ and
$\pi^2/50$, cancel it. The wall corrections are indispensable to
the existence of the flat limit.

There is also a two-dimensional exact check. For five equal weights
$\mu_i=2/5$, put $T=\sum_{i<j}\delta_{\{i,j\}}$. On
$\MM_{0,5}$ the boundary and cotangent relations give
$\sum_i\psi_i=\tfrac32T$. Each of the ten boundary curves has
self-intersection $-1$; distinct curves meet once precisely when
their index pairs are disjoint. There are fifteen such unordered
intersections, so $\int T^2=-10+2\cdot15=20$.
Thus $B_\mu=T/10$, $\Xi_\mu=T/20$, and
\[
 V_{\mathrm{CH}}=\frac85\pi^2,\qquad
 V_{\WP}(A)=\frac25\pi^2A^2+\frac15\pi A^3+\frac1{40}A^4,
 \qquad 0<A<\pi.
\]
Here $\varepsilon<1/6$ is the exact final-chamber condition.
The boundary intersections are described in
\cite[Example 2.2 and Section 5]{Kojima}; the divisor identities are
also available from~\cite[Section 7.3 and Appendix A]{KoziarzNguyen}.

The leading class $B_\mu$ and volume $V_{\mathrm{CH}}(\mu)$ vary
continuously with the weights. The higher coefficients of the
\emph{final-chamber} expansion need not. For example, let
\[
 \mu(\delta)=\left(\tfrac12-\delta,\tfrac12-\delta,
                    \tfrac12+\delta,\tfrac12+\delta\right),
 \qquad 0<\delta<\tfrac12.
\]
The only light pair is $\{1,2\}$, and
\[
 \int\Xi_{\mu(\delta)}=\frac\delta2\longrightarrow0,
 \qquad
 \int\Xi_{(1/2)^4}=-\frac18.
\]
There is no contradiction: the final chamber has
$\varepsilon<2\delta/(1+2\delta)$ and collapses as $\delta\to0$.
The full piecewise WP volume remains continuous. The small-$A$
expansion must therefore be taken at fixed weights; uniformity
across balanced collision walls requires a separate analysis.

\subsection{Why the complex hyperbolic volumes look different}

Thurston proves finite volume and connects particular examples to
counting triangulations~\cite{ThurstonShapes}. A general explicit
volume formula is obtained by McMullen's stratified Gauss--Bonnet
calculation~\cite[Theorem 8.1]{McMullenGB}. In the present labeled,
curvature-$-1$ convention it is
\begin{equation}
 V_{\mathrm{CH}}(\mu)=\frac{(-4\pi)^d}{(d+1)!}
 \sum_{\substack{\mathcal P\text{ a partition of }\{1,\ldots,n\}\\
                  |\mathcal P|\ge3}}
 (-1)^{|\mathcal P|+1}(|\mathcal P|-3)!
 \prod_{B\in\mathcal P}(1-\mu_B)_+^{|B|-1}.
 \label{eq:McMullenpartitionvolume}
\end{equation}
Here $(x)_+=\max(x,0)$. Each set partition is counted once, with
no division by permutations of equal weights; singleton factors are
$1$. Only blocks of weight less than one contribute. For the five
equal weights above, the contributions with respectively five, four,
and three blocks are $2$, $-2$, and $3/5$, giving
$V_{\mathrm{CH}}=8\pi^2/5$ independently of the intersection
calculation. This finite example checks the constant and the labeled
partition convention; the general formula is McMullen's theorem.

The blocks encode collisions. Their solid-angle factors weight
Euler characteristics of the corresponding strata. Balanced subsets
give cusps; the formula includes them by the finite-volume extension
in the proof of that theorem. Koziarz--Nguyen's
intersection formula~\eqref{eq:CHintersection} computes the same
quantity by a different route. Equation~\eqref{eq:WPCHvolumeasymptotic}
recovers it from the normalized small-area limit
$V_{\mathrm{CH}}(\mu)=2^d\lim_{A\downarrow0}A^{-d}V_{\WP}(A)$. Thus the
Gauss--Bonnet, period, and intersection calculations refer to the
same limiting geometry, with their normalization made explicit.


